\documentclass[UTF8,fontset=none]{ctexart}
\usepackage[a4paper,margin=2.5cm]{geometry}
\usepackage{graphicx}
\usepackage{tikz}
\usepackage{subcaption}
\usepackage{float}
\usepackage{xcolor}
\setCJKmainfont{Noto Serif CJK SC}
\setCJKsansfont{Noto Sans CJK SC}
\setCJKmonofont{Noto Sans Mono CJK SC}

\graphicspath{{/home/zcy/workspace/records/zjuthesis/body/graduate/paper1/figures/}}
\pagestyle{empty}

\begin{document}

\begin{figure}[H]
  \centering
  \begin{tikzpicture}
    \node [above right, inner sep=0] (image) at (0, 0) {
      \includegraphics[width=\textwidth]{helix_vs_penalty/render.png}
    };
    \begin{scope}[x={(image.south east)},y={(image.north west)}]
      \node [inner sep=0] at (0.7, 0.8) {重力$g$};
      \node [anchor=west,inner sep=0] at (0.7, 0.43) {\small{罚函数法-权重1e2}};
      \node [anchor=west,inner sep=0] at (0.7, 0.32) {\small{罚函数法-权重1e9}};
      \node [anchor=west,inner sep=0] at (0.7, 0.21) {\small{KKT方法}};
      \node [anchor=west,inner sep=0] at (0.7, 0.1) {\small{本文方法}};
    \end{scope}
  \end{tikzpicture}
  \caption{\textbf{重力作用下的螺旋线。} 模拟由 200 个线段构成的螺旋杆，每个时间步最多使用 5 次牛顿迭代。}
  \label{fig:visual_damping_penalty}
\end{figure}

\clearpage

\begin{figure}[H]
  \centering
  \begin{subfigure}[b]{\textwidth}
    \centering
    \includegraphics[width=\textwidth]{sheet/pdf/RedMax.pdf}
    \caption{扩展 RedMax 最大稳定时间步估计；紫线与橙线几乎重合。}
  \end{subfigure}

  \begin{subfigure}[b]{\textwidth}
    \centering
    \begin{tikzpicture}
      \node [above right, inner sep=0] (image) at (0, 0) {
        \includegraphics[width=\textwidth]{redMax/render2.png}
      };
      \node [anchor=east,inner sep=0] at (1.2, 3.5) {重力$g$};
      \node [anchor=south,inner sep=0] at (3.2, 5.6) {第$0$帧};
      \node [anchor=south,inner sep=0] at (8.8, 5.6) {第$1$帧};
      \node [anchor=south,inner sep=0] at (13.8, 5.6) {第$2$帧};
      \node [anchor=south west,inner sep=0] at (1.2, 2.3) {\small{RedMax($\Delta t=0.6875\mathrm{s}$)}};
      \node [anchor=south west,inner sep=0] at (1.2, 1.7) {\small{RedMax($\Delta t=0.71875\mathrm{s}$)}};
      \node [anchor=south west,inner sep=0] at (1.2, 1.1) {\small{本文方法($\Delta t=0.71875\mathrm{s}$)}};
      \node [anchor=south west,inner sep=0] at (1.2, 0.5) {\small{本文方法($\Delta t=1.3\mathrm{s}$)}};
    \end{tikzpicture}
    \caption{扩展 RedMax 与本文方法在重力螺旋场景下对比。}
  \end{subfigure}

  \caption{\textbf{与扩展 RedMax 的稳定性比较。}
  对每个重力加速度 $g$，二分法最后一步给出时间步区间（a），其起点（紫线）对应稳定仿真，终点（橙线）对应不稳定仿真。随 $g$ 增大，RedMax 稳定域显著收缩；图（b）展示了（a）中首个采样点对应的可视化结果。}
  \label{fig:timestep_limit}
\end{figure}

\end{document}
